\ifdefined\gkver\else\def\gkver{student}\fi
\documentclass[\gkver]{gaokaozhenti}
\newcommand{\ccnum}[1]{{\fontspec{Noto Serif CJK SC}#1}}
\begin{document}

\chapter{2011年安徽理综}
%% paperType: 理综物理部分

\begin{enumerate}
\item
%% number: 14
%% paperName: 2011年安徽理综第 14 题 6 分
%% typeId: 1
%% type: 单选题
%% chapter: 运动和力
%% point: 牛顿第二定律
%% method: 无
%% score: 6
%% degree: 450
%% duplicateId: 0
%% body:
一质量为 $m$ 的物块恰好静止在倾角为 $\theta$ 斜面上。现对物块施加一个竖直向下的恒力 $F$，如图所示。则物块\xzanswer{A}

\onepicture[w=4.8cm]{figs/14.png}

\fourchoices[answer=A]
{仍处于静止状态}
{沿斜面加速下滑}
{受到的摩擦力不变}
{受到的合外力增大}

%% answer: A
%% memo:
\memoanswer{
由于质量为 $m$ 的物块恰好静止在倾角为 $\theta$ 的斜面上，说明斜面对物块的作用力与物块的重力平衡，斜面与物块的动摩擦因数 $\mu=\tan\theta$。对物块施加一个竖直向下的恒力 $F$，使得合力仍然为零，故物块仍处于静止状态，A 正确，B、D 错误。摩擦力由 $mg\sin\theta$ 增大到 $(F+mg)\sin\theta$，C 错误。
}

\item
%% number: 15
%% paperName: 2011年安徽理综第 15 题 6 分
%% typeId: 1
%% type: 单选题
%% chapter: 光学
%% point: 光的折射
%% method: 无
%% score: 6
%% degree: 650
%% duplicateId: 0
%% body:
实验表明，可见光通过三棱镜时各色光的折射率 $n$ 随着波长 $\lambda$ 的变化符合科西经验公式：$n=A+\frac{B}{\lambda^{2}}+\frac{C}{\lambda^{4}}$，其中 A、B、C 是正的常量。太阳光进入三棱镜后发生色散的情形如下图所示。则\xzanswer{D}

\onepicture[w=8cm]{figs/15.png}

\fourchoices[answer=D]
{屏上 c 处是紫光}
{屏上 d 处是红光}
{屏上 b 处是紫光}
{屏上 a 处是红光}

%% answer: D
%% memo:
\memoanswer{
白色光经过三棱镜后产生色散现象，在光屏由上至下（a、b、c、d）依次为红、橙、黄、绿、蓝、靛、紫。屏上 a 处为红光，屏上 d 处是紫光，D 正确。
}

\item
%% number: 16
%% paperName: 2011年安徽理综第 16 题 6 分
%% typeId: 1
%% type: 单选题
%% chapter: 直线运动
%% point: 时间中点速度
%% method: 无
%% score: 6
%% degree: 450
%% duplicateId: 0
%% body:
一物体作匀加速直线运动，通过一段位移 $\Delta x$ 所用的时间为 $t_{1}$，紧接着通过下一段位移 $\Delta x$ 所用时间为 $t_{2}$。则物体运动的加速度为\xzanswer{A}

\fourchoices[answer=A]
{$\frac{2\Delta x(t_{2}-t_{1})}{t_{1}t_{2}(t_{1}+t_{2})}$}
{$\frac{\Delta x(t_{2}-t_{1})}{t_{1}t_{2}(t_{1}+t_{2})}$}
{$\frac{2\Delta x(t_{1}+t_{2})}{t_{1}t_{2}(t_{1}-t_{2})}$}
{$\frac{\Delta x(t_{1}+t_{2})}{t_{1}t_{2}(t_{1}-t_{2})}$}

%% answer: A
%% memo:
\memoanswer{
物体作匀加速直线运动在前一段 $\Delta x$ 所用的时间为 $t_{1}$，平均速度为 $\overline{v_{1}}=\frac{\Delta x}{t_{1}}$，即为 $\frac{t_{1}}{2}$ 时刻的瞬时速度；物体在后一段 $\Delta x$ 所用的时间为 $t_{2}$，平均速度为 $\overline{v_{2}}=\frac{\Delta x}{t_{2}}$，即为 $\frac{t_{2}}{2}$ 时刻的瞬时速度。速度由 $\overline{v_{1}}$ 变化到 $\overline{v_{2}}$ 的时间为 $\Delta t=\frac{t_{1}+t_{2}}{2}$，所以加速度 $a=\frac{\overline{v_{2}}-\overline{v_{1}}}{\Delta t}=\frac{2\Delta x(t_{1}-t_{2})}{t_{1}t_{2}(t_{1}+t_{2})}$，A 正确。
}

\item
%% number: 17
%% paperName: 2011年安徽理综第 17 题 6 分
%% typeId: 1
%% type: 单选题
%% chapter: 曲线运动
%% point: 圆周运动的应用
%% method: 无
%% score: 6
%% degree: 450
%% duplicateId: 0
%% body:
一般的曲线运动可以分成很多小段，每小段都可以看成圆周运动的一部分，即把整条曲线用一系列不同半径的小圆弧来代替。如图（a）所示，曲线上的 A 点的曲率圆定义为：通过 A 点和曲线上紧邻 A 点两侧的两点作一圆，在极限情况下，这个圆就叫做 A 点的曲率圆，其半径 $\rho$ 叫做 A 点的曲率半径。现将一物体沿与水平面成 $\alpha$ 角的方向以速度 $v_{0}$ 抛出，如图（b）所示。则在其轨迹最高点 P 处的曲率半径是\xzanswer{C}

\onepicture[w=6cm]{\includesvg{figs/17}}

\fourchoices[answer=C]
{$\frac{v_{0}^{2}}{g}$}
{$\frac{v_{0}^{2}\sin^{2}\alpha}{g}$}
{$\frac{v_{0}^{2}\cos^{2}\alpha}{g}$}
{$\frac{v_{0}^{2}\cos^{2}\alpha}{g\sin\alpha}$}

%% answer: C
%% memo:
\memoanswer{
物体在其轨迹最高点 P 处只有水平速度，其水平速度大小为 $v_{0}\cos\alpha$，根据牛顿第二定律得 $mg=m\frac{(v_{0}\cos\alpha)^{2}}{\rho}$，所以在其轨迹最高点 P 处的曲率半径是 $\rho=\frac{v_{0}^{2}\cos^{2}\alpha}{g}$，C 正确。
}

\item
%% number: 18
%% paperName: 2011年安徽理综第 18 题 6 分
%% typeId: 1
%% type: 单选题
%% chapter: 电场
%% point: 带电粒子的偏转
%% method: 无
%% score: 6
%% degree: 450
%% duplicateId: 0
%% body:
图（a）为示波管的原理图。如果在电极 YY$'$ 之间所加的电压按图（b）所示的规律变化，在电极 XX$'$ 之间所加的电压按图（c）所示的规律变化，则在荧光屏上会看到的图形是\xzanswer{B}

\onepicture[w=11cm]{figs/18.png}

\fourchoices[answer=B, ispicture=true, h=2.4cm]
{figs/18a.png}
{figs/18b.png}
{figs/18c.png}
{figs/18d.png}

%% answer: B
%% memo:
\memoanswer{
由于电极 XX$'$ 加的是扫描电压，电极 YY$'$ 之间所加的电压是信号电压，所以荧光屏上会看到的图形是 B，答案 B 正确。
}

\item
%% number: 19
%% paperName: 2011年安徽理综第 19 题 6 分
%% typeId: 1
%% type: 单选题
%% chapter: 电磁感应
%% point: 切割磁感线电动势
%% method: 无
%% score: 6
%% degree: 450
%% duplicateId: 0
%% body:
如图所示的区域内有垂直于纸面的匀强磁场，磁感应强度为 $B$。电阻为 $R$、半径为 $L$、圆心角为 $45^\circ$ 的扇形闭合导线框绕垂直于纸面的 O 轴以角速度 $\omega$ 匀速转动（O 轴位于磁场边界）。则线框内产生的感应电流的有效值为\xzanswer{D}

\onepicture[w=4.2cm]{figs/19.png}

\fourchoices[answer=D]
{$\frac{BL^{2}\omega}{2R}$}
{$\frac{\sqrt{2}BL^{2}\omega}{2R}$}
{$\frac{\sqrt{2}BL^{2}\omega}{4R}$}
{$\frac{BL^{2}\omega}{4R}$}

%% answer: D
%% memo:
\memoanswer{
交流电流的有效值是根据电流的热效应得出的，线框转动周期为 $T$，而线框转动一周只有 $T/4$ 的时间内有感应电流，则有 $\left(\frac{BL\omega\frac{L}{2}}{R}\right)^{2}R\frac{T}{4}=I^{2}RT$，所以 $I=\frac{BL^{2}\omega}{4R}$。D 正确。
}

\item
%% number: 20
%% paperName: 2011年安徽理综第 20 题 6 分
%% typeId: 1
%% type: 单选题
%% chapter: 电场
%% point: 带电粒子的直线运动
%% method: 无
%% score: 6
%% degree: 450
%% duplicateId: 0
%% body:
如图（a）所示，两平行正对的金属板 A、B 间加有如图（b）所示的交变电压，一重力可忽略不计的带正电粒子被固定在两板的正中间 P 处。若在 $t_{0}$ 时刻释放该粒子，粒子会时而向 A 板运动，时而向 B 板运动，并最终打在 A 板上。则 $t_{0}$ 可能属于的时间段是\xzanswer{B}

\twopicture[numA=图（a）, numB=图（b）, labelprefix={}, wA=3.8cm, wB=5.2cm]
{figs/20a.png}
{figs/20b.png}

\fourchoices[answer=B]
{$0<t_{0}<T/4$}
{$T/2<t_{0}<3T/4$}
{$3T/4<t_{0}<T$}
{$T<t_{0}<9T/8$}

%% answer: B
%% memo:
\memoanswer{
若 $0<t_{0}<\frac{T}{4}$，带正电粒子先加速向 B 板运动、再减速运动至零；然后再反方向加速运动、减速运动至零；如此反复运动，每次向右运动的距离大于向左运动的距离，最终打在 B 板上，所以 A 错误。若 $\frac{T}{2}<t_{0}<\frac{3T}{4}$，带正电粒子先加速向 A 板运动、再减速运动至零；然后再反方向加速运动、减速运动至零；如此反复运动，每次向左运动的距离大于向右运动的距离，最终打在 A 板上，所以 B 正确。若 $\frac{3T}{4}<t_{0}<T$，带正电粒子先加速向 A 板运动、再减速运动至零；然后再反方向加速运动、减速运动至零；如此反复运动，每次向左运动的距离小于向右运动的距离，最终打在 B 板上，所以 C 错误。若 $T<t_{0}<\frac{9T}{8}$，带正电粒子先加速向 B 板运动、再减速运动至零；然后再反方向加速运动、减速运动至零；如此反复运动，每次向右运动的距离大于向左运动的距离，最终打在 B 板上，所以 D 错误。
}

\item
%% number: 21
%% paperName: 2011年安徽理综第 21 题 12 分
%% typeId: 4
%% type: 实验题
%% chapter: 电路
%% point: 多用表的使用
%% method: 无
%% score: 12
%% degree: 450
%% duplicateId: 0
%% body:
\begin{enumerate}
	\item
	\textbf{Ⅰ．}为了测量某一弹簧的劲度系数，将该弹簧竖直悬挂起来，在自由端挂上不同质量的砝码。实验测出了砝码的质量 $m$ 与弹簧长度 $l$ 的相应数据，七对应点已在图上标出。（$g=9.8\Umsq$）
	\begin{enumerate}
		\item
		作出 $m-l$ 的关系图线；
		\item
		弹簧的劲度系数为 \tkanswer[2cm]{0.248$\sim$0.262} $\UN/\Um$。
	\end{enumerate}

	\onepicture[align=right, w=6.5cm]{figs/21.png}

	\item
	\textbf{Ⅱ．}
	\begin{enumerate}
		\item
		某同学使用多用电表粗略测量一定值电阻的阻值，先把选择开关旋到“$\times1k$”挡位，测量时针偏转如图（a）所示。请你简述接下来的测量过程：
		\begin{steps}
			\item
			\tkanswer[6cm]{断开待测电阻，将选择开关旋到“$\times100$”档}；
			\item
			\tkanswer[6cm]{将两表笔短接，调整“欧姆调零旋钮”，使指针指向“$0\UO$”}；
			\item
			\tkanswer[6cm]{再接入待测电阻，将指针示数$\times100$，即为待测电阻阻值}；
			\item
			测量结束后，将选择开关旋到“OFF”挡。
		\end{steps}
		\item
		接下来采用“伏安法”较准确地测量该电阻的阻值，所用实验器材如图（b）所示。其中电压表内阻约为 $5\UkO$，电流表内阻约为 $5\UO$。图中部分电路已经连接好，请完成实验电路的连接。
		\item
		图（c）是一个多量程多用电表的简化电路图，测量电流、电压和电阻各有两个量程。当转换开关 S 旋到位置 3 时，可用来测量 \tkanswer[1.4cm]{电阻}；当 S 旋到位置 \tkanswer[1.6cm]{1、2} 时，可用来测量电流，其中 S 旋到位置 \tkanswer[1.2cm]{1} 时量程较大。
	\end{enumerate}

	\threepicture[align=right, numA=图（a）, numB=图（b）, numC=图（c）, labelprefix={}, wA=3.2cm, wB=4.6cm, wC=4.6cm]
	{figs/21a.png}
	{figs/21b.png}
	{figs/21c.png}
\end{enumerate}

%% answer: 
\jdanswer{
\begin{enumerate}
	\item
	Ⅰ．\begin{enumerate}
		\item
		图线如图所示（略）；
		\item
		$0.248\sim0.262\UN/\Um$。
	\end{enumerate}
	\item
	Ⅱ．\begin{enumerate}
		\item
		①断开待测电阻，将选择开关旋到“$\times100$”档；②将两表笔短接，调整“欧姆调零旋钮”，使指针指向“$0\UO$”；③再接入待测电阻，将指针示数$\times100$，即为待测电阻阻值。
		\item
		如图（b）所示。
		\item
		电阻，1、2，1。
	\end{enumerate}
\end{enumerate}
}
%% memo:
\memoanswer{
Ⅰ．（1）如图所示。

\onepicture[w=6.5cm]{figs/21_an1.png}

（2）$0.248\sim0.262\UN/\Um$。

Ⅱ．（1）①断开待测电阻，将选择开关旋到“$\times100$”档；②将两表笔短接，调整“欧姆调零旋钮”，使指针指向“$0\UO$”；③再接入待测电阻，将指针示数$\times100$，即为待测电阻阻值。

（2）如图所示。

\onepicture[w=5.4cm]{figs/21_an2.png}

（3）电阻，1、2，1。
}

\item
%% number: 22
%% paperName: 2011年安徽理综第 22 题 14 分
%% typeId: 6
%% type: 计算题
%% chapter: 曲线运动
%% point: 天体运动
%% method: 无
%% score: 14
%% degree: 450
%% duplicateId: 0
%% body:
\begin{enumerate}
	\item
	开普勒行星运动第三定律指出：行星绕太阳运动的椭圆轨道的半长轴 $a$ 的三次方与它的公转周期 $T$ 的二次方成正比，即 $\frac{a^{3}}{T^{2}}=k$，$k$ 是一个对所有行星都相同的常量。将行星绕太阳的运动按圆周运动处理，请你推导出太阳系中该常量 $k$ 的表达式。已知引力常量为 $G$，太阳的质量为 $M_{\text{太}}$。
	\item
	开普勒定律不仅适用于太阳系，它对一切具有中心天体的引力系统（如地月系统）都成立。经测定月地距离为 $3.84\times10^{8}\Um$，月球绕地球运动的周期为 $2.36\times10^{6}\Us$，试计算地球的质量 $M_{\text{地}}$。（$G=6.67\times10^{-11}\UNmqkgq$，结果保留一位有效数字）
\end{enumerate}

%% answer: 
\jdanswer{
\begin{enumerate}
	\item
	$k=\frac{GM_{\text{太}}}{4\pi^{2}}$
	\item
	$M_{\text{地}}=6\times10^{24}\Ukg$（$M_{\text{地}}=5\times10^{24}\Ukg$ 也算对）
\end{enumerate}
}
%% memo:
\memoanswer{
（1）因行星绕太阳作匀速圆周运动，于是轨道的半长轴 $a$ 即为轨道半径 $r$。根据万有引力定律和牛顿第二定律有

$G\frac{M_{\text{太}}m_{\text{行}}}{r^{2}}=m_{\text{行}}\frac{4\pi^{2}}{T^{2}}r$ ①

于是有

$\frac{r^{3}}{T^{2}}=\frac{GM_{\text{太}}}{4\pi^{2}}$ ②

即

$k=\frac{GM_{\text{太}}}{4\pi^{2}}$ ③

（2）在月地系统中，设月球绕地球运动的轨道半径为 $R$，周期为 $T$，由②式可得

$\frac{R^{3}}{T^{2}}=\frac{GM_{\text{地}}}{4\pi^{2}}$ ④

解得

$M_{\text{地}}=6\times10^{24}\Ukg$ ⑤

（$M_{\text{地}}=5\times10^{24}\Ukg$ 也算对）
}

\item
%% number: 23
%% paperName: 2011年安徽理综第 23 题 16 分
%% typeId: 6
%% type: 计算题
%% chapter: 磁场
%% point: 带电粒子在磁场中的运动
%% method: 无
%% score: 16
%% degree: 450
%% duplicateId: 0
%% body:
如图所示，在以坐标原点 O 为圆心、半径为 $R$ 的半圆形区域内，有相互垂直的匀强电场和匀强磁场，磁感应强度为 $B$，磁场方向垂直于 $xOy$ 平面向里。一带正电的粒子（不计重力）从 O 点沿 $y$ 轴正方向以某一速度射入，带电粒子恰好做匀速直线运动，经 $t_{0}$ 时间从 P 点射出。
\begin{enumerate}
	\item
	求电场强度的大小和方向。
	\item
	若仅撤去磁场，带电粒子仍从 O 点以相同的速度射入，经 $\frac{t_{0}}{2}$ 时间恰从半圆形区域的边界射出。求粒子运动加速度的大小。
	\item
	若仅撤去电场，带电粒子仍从 O 点射入，且速度为原来的 4 倍，求粒子在磁场中运动的时间。
\end{enumerate}
\onepicture[align=right, w=5.5cm]{figs/23.png}
%% answer: 
\jdanswer{
\begin{enumerate}
	\item
	$E=\frac{BR}{t_{0}}$
	\item
	$a=\frac{4\sqrt{3}R}{t_{0}^{2}}$
	\item
	$t=\frac{\sqrt{3}\pi}{18}t_{0}$
\end{enumerate}
}
%% memo:
\memoanswer{
（1）设带电粒子的质量为 $m$，电荷量为 $q$，初速度为 $v$，电场强度为 $E$。可判断出粒子受到的洛伦兹力沿 $x$ 轴负方向，于是可知电场强度沿 $x$ 轴正方向

且有

$qE=qvB$ ①

又

$R=vt_{0}$ ②

则

$E=\frac{BR}{t_{0}}$ ③

（2）仅有电场时，带电粒子在匀强电场中作类平抛运动

在 $y$ 方向位移

$y=v\frac{t_{0}}{2}$ ④

由②④式得

$y=\frac{R}{2}$ ⑤

设在水平方向位移为 $x$，因射出位置在半圆形区域边界上，于是

$x=\frac{\sqrt{3}}{2}R$

又有

$x=\frac{1}{2}a\left(\frac{t_{0}}{2}\right)^{2}$ ⑥

得

$a=\frac{4\sqrt{3}R}{t_{0}^{2}}$ ⑦

（3）仅有磁场时，入射速度 $v'=4v$，带电粒子在匀强磁场中作匀速圆周运动，设轨道半径为 $r$，由牛顿第二定律有

$qv'B=m\frac{v'^{2}}{r}$ ⑧

又

$qE=ma$ ⑨

由⑦⑧⑨式得

$r=\frac{\sqrt{3}}{3}R$ \ccnum{⑩}

由几何关系

$\sin\alpha=\frac{R}{2r}$ \ccnum{⑪}

即

$\sin\alpha=\frac{\sqrt{3}}{2}$，$\alpha=\frac{\pi}{3}$ \ccnum{⑫}

带电粒子在磁场中运动周期

$T=\frac{2\pi m}{qB}$

则带电粒子在磁场中运动时间

$t=\frac{2\alpha}{2\pi}T$

所以

$t=\frac{\sqrt{3}\pi}{18}t_{0}$ \ccnum{⑬}

\onepicture[w=5cm]{figs/23_an.png}
}

\item
%% number: 24
%% paperName: 2011年安徽理综第 24 题 20 分
%% typeId: 6
%% type: 计算题
%% chapter: 功和能
%% point: 机械能守恒定律
%% method: 微元
%% score: 20
%% degree: 250
%% duplicateId: 0
%% body:
如图所示，质量 $M=2\Ukg$ 的滑块套在光滑的水平轨道上，质量 $m=1\Ukg$ 的小球通过长 $L=0.5\Um$ 的轻质细杆与滑块上的光滑轴 O 连接，小球和轻杆可在竖直平面内绕 O 轴自由转动，开始轻杆处于水平状态，现给小球一个竖直向上的初速度 $v_{0}=4\Ums$，$g$ 取 $10\Umsq$。
\begin{enumerate}
	\item
	若锁定滑块，试求小球通过最高点 P 时对轻杆的作用力大小和方向。
	\item
	若解除对滑块的锁定，试求小球通过最高点时的速度大小。
	\item
	在满足（2）的条件下，试求小球击中滑块右侧轨道位置点与小球起始位置点间的距离。
\end{enumerate}
\onepicture[align=right, w=6.5cm]{figs/24.png}
%% answer: 
\jdanswer{
\begin{enumerate}
	\item
	$F=2\UN$，方向竖直向上。
	\item
	$v_{2}=2\Ums$
	\item
	$s_{1}=\frac{2}{3}\Um$
\end{enumerate}
}
%% memo:
\memoanswer{
（1）设小球能通过最高点，且此时的速度为 $v_{1}$。在上升过程中，因只有重力做功，小球的机械能守恒。则

$\frac{1}{2}mv_{1}^{2}+mgL=\frac{1}{2}mv_{0}^{2}$ ①

$v_{1}=\sqrt{6}\Ums$ ②

设小球到达最高点时，轻杆对小球的作用力为 $F$，方向向下，则

$mg+F=m\frac{v_{1}^{2}}{L}$ ③

由②③式，得

$F=2\UN$ ④

由牛顿第三定律可知，小球对轻杆的作用力大小为 $2\UN$，方向竖直向上。

（2）解除锁定后，设小球通过最高点时的速度为 $v_{2}$，此时滑块的速度为 $V$。在上升过程中，因系统在水平方向上不受外力作用，水平方向的动量守恒。以水平向右的方向为正方向，有

$mv_{2}+MV=0$ ⑤

在上升过程中，因只有重力做功，系统的机械能守恒，则

$\frac{1}{2}mv_{2}^{2}+\frac{1}{2}MV^{2}+mgL=\frac{1}{2}mv_{0}^{2}$ ⑥

由⑤⑥式，得

$v_{2}=2\Ums$ ⑦

（3）设小球击中滑块右侧轨道的位置点与小球起始点的距离为 $s_{1}$，滑块向左移动的距离为 $s_{2}$，任意时刻小球的水平速度大小为 $v_{3}$，滑块的速度大小为 $V'$。由系统水平方向的动量守恒，得

$mv_{3}-MV'=0$ ⑧

将⑧式两边同乘以 $\Delta t$，得

$mv_{3}\Delta t-MV'\Delta t=0$ ⑨

因⑨式对任意时刻附近的微小间隔 $\Delta t$ 都成立，累积相加后，有

$ms_{1}-Ms_{2}=0$ \ccnum{⑩}

又

$s_{1}+s_{2}=2L$ \ccnum{⑪}

由\ccnum{⑩⑪}式得

$s_{1}=\frac{2}{3}\Um$ \ccnum{⑫}
}

\end{enumerate}

\end{document}
